% from main.tex
%%%%%%%%%%  scan idx 395  |  folio 382  %%%%%%%%%%
donne par adjonction un homomorphisme~:
\[
  \bar f^{\,*}(M)\longrightarrow M
\]
et on vérifie immédiatement sur chaque fibre que c'est un isomorphisme.

\textbf{2.3.2.} Soient maintenant \(B\) une \(\mathcal O_S\)-Algèbre graduée quasi-cohérente, \(M\) un \(B\)-Module gradué, et \(X=\Proj(B)\). Pour tout ouvert affine \(U\) de \(S\), et toute section \(s\) de \(B\) sur \(U\), homogène de degré \(d>0\), on note \(M_{(s),U}\) la composante homogène de degré \(0\) du Module de fraction \(M_s\); on a donc :
\[
  M_{(s),U}=
  \bigl((M|_U)\otimes_{(B|_U)}(B|_U)_s\bigr)_0 .
\]
Si \(V\) est un ouvert affine de \(U\), alors \(M_{(s),V}\) est la restriction à \(V\) de \(M_{(s),U}\).

On considère la base d'ouverts de \(X\) formée des ouverts affines :
\[
  D_+(U,s)\simeq \Spec\bigl((B|_U)_{(s)}\bigr) .
\]
On note \(f\) le morphisme structural \(X\to S\), et \(f_{U,s}\) le morphisme affine :
\[
  f_{U,s}:D_+(U,s)\longrightarrow U .
\]
On pose alors :
\[
  \Proj_0(M)_{U,s}
  =
  \bar f_{U,s}^{\,*}\bigl(M_{(s),U}\bigr),
\]
avec les notations de 2.3.1. Nous allons montrer que les \(\Proj_0(M)_{U,s}\) se recollent quand \(U\) et \(s\) varient, définissant ainsi un \(\mathcal O_X\)-Module noté \(\Proj_0(M)\). Le recollement pour \(U\) variable et \(s\) fixe étant évident, on est ramené (cf. EGA II 2.5.2) à montrer que si \(t\) est une section de \(B\) sur \(U\), homogène de degré \(e>0\), alors la restriction de \(\Proj_0(M)_{U,s}\) à \(D_+(U,st)\) est canoniquement isomorphe à \(\Proj_0(M)_{U,st}\). On considère :
